\section{Scope, main advances, and source status}
\label{bc:scope}

This article is a continuation of the ProveIt material at the immutable
repository checkpoint
\[
 \pin .
\]
The requested manuscript directory, the seven archives in
\src{docs/incoming}, and the directly relevant preserved continuations were
inspected together. This matters: an open question in one report may have
been solved in another report that has not yet been integrated into the
canonical manuscript. The accompanying provenance file records the
identities and hashes of the sources used here. The source tree itself was
not modified.

\subsection{The questions answered}

The first precise target is the question headed \emph{Every bilinear contact
hyperplane} in \emph{Bilinear Tails and Determinants}
\cite{ProveItBilinear}. Its finite two-order endpoint identity is already
proved. The missing data are the bivariate contact polynomials on
\(p+q=-2m\), all independent order derivatives, and the changes under
translation of the cutoff. A diagonal restriction cannot recover
derivatives transverse to that restriction. Section~\ref{ct:section}
supplies these data for every \(m\ge0\), including explicit ordinary
convergent representatives. The finite endpoint law and its first
\(1/2\) contact remain attributed to the predecessor.

The second target is the higher-order, many-mass question in
\emph{Higher Quadratic Jets}, item 3, and \emph{Beta Polarization of Higher
Spectral Defects}, item 3
\cite{ProveItHigherQuadratic,ProveItHigherDefects}. The latter already gives
a complete pairwise formula for second jets. Its research question asks
for higher connected logarithmic terms and correctly normalized
coincident-mass limits. Section~\ref{sp:section} gives a single real
resolvent integral for every scalar jet and a compact integral for every
connected logarithmic product. The number of real integration variables
stays one as both the number of masses and the derivative order increase.
Section~\ref{ng:section} further extends the full-jet formula to every
even polynomial weight and to every nonpositive integer spectral point.
This settles the scalar representation and connected-product parts of the
question, and supplies an alternative explicit representation for the
higher negative-point jets. It does not assert that all those integrals reduce to finitely
many special values of a smaller function class.

\begin{center}
\small
\renewcommand{\arraystretch}{1.25}
\begin{tabular}{>{\raggedright\arraybackslash}p{0.22\textwidth}>{\raggedright\arraybackslash}p{0.31\textwidth}>{\raggedright\arraybackslash}p{0.38\textwidth}}
\toprule
Topic & Result in this article & Status relative to inspected sources\\
\midrule
Two-order tail contacts &
Finite coefficients at every resonance; all mixed jets and translations &
Completes the explicitly open bivariate contact question.\\
Ordinary boundary sums &
Polynomial and logarithmic subtractions, with exact examples &
Consequences of the completed contact algebra.\\
Full scalar spectral jets &
One digamma integral for every \(N,r\) &
Explicit all-order representation beyond the supplied second-jet formulas.\\
Even weights and negative points &
Every polynomial degree and every derivative at each nonpositive integer &
A finite Hurwitz baseline and the same single ordinary integral.\\
Connected logarithmic products &
Compact density, moments, critical odd weight, collisions &
All-order realization of the higher connected-term target.\\
Local spectral correction &
Finite residue polynomial at fixed \(N\) &
Inherited mechanism, restated with matching normalization.\\
Harmonic and polylogarithmic forms &
Beta moments, confluent expansions, height-one transform &
Derived bridge to classical coefficient and function alphabets.\\
\bottomrule
\end{tabular}
\end{center}

\subsection{Two organizing mechanisms}

The tail calculation is controlled by an \emph{upper endpoint}. Its
asymptotic coefficient at a resonant exponent gives the discrepancy between
an analytic continuation and a specified cutoff finite part. The necessary
operation is to expand in independent parameters before restricting to the
resonance. A finite-degree theorem then makes all sufficiently high
independent contact derivatives vanish. The same asymptotic coefficients
specify actual convergent subtractions, so the finite-part formulas have
ordinary-series counterparts.

The spectral calculation is controlled by a \emph{logarithmic jump across
a cut}. Subtracting the mass-zero logarithm removes the obstruction at
infinity. Cauchy's formula converts the remaining logarithmic expression
into a resolvent integral, and the lattice resolvent sums exactly to a
digamma difference. For a product of logarithmic mass increments the cut
is compact and the decay at infinity forces its first moments to vanish.
Those moment cancellations supply the polynomial subtractions for weighted
sums.

These are different mechanisms. The cutoff degree theorem is a statement
about asymptotic coefficients in the orders \(p,q\). The many-mass
interaction hierarchy is a statement about finite differences of spectral
jets with a fixed total number of factors. Neither follows by replacing
one meaning of degree with the other.

\subsection{Classical inputs and the priority boundary}

The Gamma quotient expansion, the alternating Hurwitz expansion and its
parameter derivatives, the digamma resolvent identity, and Euler's beta
integral are classical
\cite{DLMF,HuKim,WangHuKim}. Products of zeta tails have an established
multiple-zeta literature; Hoffman's paper is a relevant example
\cite{Hoffman}. The broad relation between regularized traces and local
asymptotic coefficients is likewise established
\cite{PaychaScott}. The normalized exponent-local degree theorem is already
present in the project's \emph{Arithmetic Spectral Jets} continuation
\cite{ProveItArithmeticJets}. We use it for comparison and give a direct
normalization check where it enters.

The contribution claimed here is the stated collection of precise
deductions and completed research targets relative to the inspected
project material. The article gives proofs of those statements. A targeted
search did not establish an earlier occurrence of every displayed form,
but that is not a comprehensive historical-priority result. In particular,
we do not relabel a classical beta integral, Stirling-number identity, or
height-one multiple-zeta reduction as a newly discovered constant.

\subsection{Conventions}

\(\zeta^{(r)}(s,a)\) always means an \emph{order} derivative
\(\partial_s^r\zeta(s,a)\). Superscripts on \(\psi\) mean argument
derivatives. The generalized Stieltjes convention is
\[
 \zeta(1+u,a)=\frac1u+
   \sum_{j\ge0}\frac{(-1)^j}{j!}\gamma_j(a)u^j,
 \qquad \gamma_0(a)=-\psi(a).
\]
The generalized harmonic numbers are
\[
 H_n^{(j)}=\sum_{k=1}^{n}k^{-j},\qquad H_0^{(j)}=0,
 \qquad H_n=H_n^{(1)}.
\]
A jet of order \(r\) is the actual derivative, with its factorial included;
\([u^r]\) instead denotes an ordinary Taylor coefficient.

Every finite part is tied to its stated coordinate or independent
regulator. For the tail problem, the coordinate is \(X=N+c\) and the
retained term is \(X^0(\log X)^0\). For the alternating endpoint a parity
subsequence is also specified. All integrals labeled ordinary are
absolutely convergent improper integrals unless conditional convergence is
explicitly stated. Signed logarithmic densities are allowed; positivity
is not assumed or required.
